\subsection{滤代数中的指数与对数}

设 \(\mathbf{A}\) 是一个含幺结合代数，在实滤子 \((\mathbf{A}_{\alpha})\) 下是豪斯多夫且完备的。记 \(m = A_0^+ = \bigcup_{\alpha > 0} A_\alpha\)。

对于 \(x \in \mathfrak{m}\)，族 \((x^n / n!)_{n \in \mathbf{N}}\) 是可求和的。记

\[
e ^ {x} = \exp x = \sum_ {n \geqslant 0} x ^ {n} / n!.\tag{1}
\]

于是 \(\exp (x)\in 1 + \mathfrak{m}\)，映射 \(\exp \colon \mathfrak{m}\to 1 + \mathfrak{m}\) 称为 A 的指数映射。

对于所有 \(y \in 1 + \mathfrak{m}\)，族 \(((-1)^{n-1}(y - 1)^n / n)_{n \geqslant 1}\) 是可求和的。记

\[
\log y = \sum_ {n \geqslant 1} (- 1) ^ {n - 1} (y - 1) ^ {n} / n.\tag{2}
\]

于是 \(\log y \in m\)，映射 \(\log: 1 + m \to m\) 称为 A 的对数映射。

命题 1。指数映射是从 \(\mathfrak{m}\) 到 \(1 + \mathfrak{m}\) 的同胚，而对数映射是其逆同胚。

对于 \(x \in \mathrm{A}_{\alpha}, \frac{x^n}{n!} \in \mathrm{A}_{n\alpha}\)，有 \(\mathrm{A}_{\alpha}\)。因此，定义指数的级数在每个 \(\alpha > 0\) 的集合 \(\mathrm{A}_{\alpha}\) 上一致收敛；由于 \(m\) 在 m 中是开集且 \(m = \bigcup_{\alpha > 0} \mathrm{A}_{\alpha}\)，指数映射连续。类似地可证明对数映射连续。

设 \(e\) 和 \(l\) 是没有常数项的形式幂级数

\[
e (\mathrm{X}) = \sum_ {n \geqslant 1} \frac {\mathrm{X} ^ {n}}{n !}, \quad l (\mathrm{X}) = \sum_ {n \geqslant 1} (- 1) ^ {n - 1} \mathrm{X} ^ {n} / n.
\]

我们知道（《代数》，第 IV 章，§ 6，第 9 号），在 \(e(l(\mathbf{X})) = l(e(\mathbf{X})) = \mathbf{X}\) 上有 \(\hat{\mathrm{A}}(\{\mathbf{X}\}) = \mathrm{K}[[\mathbf{X}]]\)。由代入（§ 5，第 1 号），得到

\[
e (l (x)) = l (e (x)) = x
\]

对于 \(x\in \mathfrak{m}\)；由于

\[
\exp x = e (x) + 1, \quad \log (1 + x) = l (x)
\]

于是立即得到

\[
\log \exp x = x, \quad \exp \log (1 + x) = 1 + x
\]

对于 \(x\) 中的 \(\mathfrak{m}\)，从而命题得证。

注。（1）若 \(x \in \mathfrak{m}\)、\(y \in \mathfrak{m}\) 且 \(x\) 与 \(y\) 可交换，则

\[
\exp (x + y) = \exp (x) \exp (y),
\]

因为族 \(\left(\frac{x^i}{i!}\cdot \frac{y^j}{j!}\right)_{i,j\in \mathbf{N}}\) 是可求和的（参见《代数》，第 IV 章，§ 6，第 9 号，命题 11）。

\begin{enumerate}
\def\labelenumi{(\arabic{enumi})}
\setcounter{enumi}{1}
\item
  由于级数 \(e\) 和 \(l\) 不含常数项且 \(A_{\alpha}\) 是 \(A\) 的闭理想，有 \(\exp A_{\alpha} \subset 1 + A_{\alpha}\) 和 \(\log (1 + A_{\alpha}) \subset A_{\alpha}\)，从而当 \(\exp A_{\alpha} = 1 + A_{\alpha}\) 时，\(\log (1 + A_{\alpha}) = A_{\alpha}\) 且 \(\alpha > 0\)。
\item
  设 B 是一个完备豪斯多夫滤含幺结合代数，且 \(n = \bigcup_{\alpha > 0} B_{\alpha}\)。设 f 是从 A 到 B 的连续含幺同态，满足 \(f(\mathfrak{m}) \subset \mathfrak{n}\)。则对于 \(f(\exp x) = \exp f(x)\)，有 \(x \in \mathfrak{m}\)；对于 \(f(\log y) = \log f(y)\)，有 \(y \in 1 + \mathfrak{m}\)；例如我们证明第一个公式：
\end{enumerate}

\[
f (\exp x) = \sum_ {n \geqslant 0} f (x ^ {n}) / n! = \sum_ {n \geqslant 0} f (x) ^ {n} / n! = \exp f (x).
\]

\begin{enumerate}
\def\labelenumi{(\arabic{enumi})}
\setcounter{enumi}{3}
\tightlist
\item
  设 E 是一个含幺结合代数。若 \(a\) 是 E 的幂零元素，则族 \(\left(\frac{a^n}{n!}\right)_{n\in \mathbf{N}}\) 具有有限支集，记 \(\exp a = \sum_{n\geqslant 0}a^n /n!\)。若 \(b\) 是幂零的，则元素 \(b - 1\) 称为幂单元；于是记
\end{enumerate}

\[
\log b = \sum_ {n \geqslant 1} (- 1) ^ {n - 1} (b - 1) ^ {n} / n.
\]

由关系 \(e(l(\mathbf{X})) = l(e(\mathbf{X})) = \mathbf{X}\) 可推出，映射 \(a \mapsto \exp a\) 是从 E 的幂零元素集到 E 的幂单元元素集的双射，而 \(b \mapsto \log b\) 是其逆映射。

